\chapter{Integrals and the Fundamental Theorem}\label{ch:integrals}

An integral is a rational-box algorithm, not an ambient area supplied by
completeness.  Its motivating calculation is the area below the arctangent
kernel.  The same design handles every admitted special function.

\begin{verbatim}
partition -> certified value ranges on cells -> add rectangles
          -> nested boxes -> shrink-width certificate
\end{verbatim}

The point is to relate all such computations, crude or fast. Lebesgue measure
and unrestricted measurable functions are outside this foundation.

\section{Rectangles}\label{sec:integrals-ftc-contract}

For a partition \(a=x_0<\cdots<x_m=b\), cell ranges \([L_k,U_k]\) give
\[
 \sum_k (x_{k+1}-x_k)L_k\ \leq\ \int_a^b f\ \leq
 \sum_k (x_{k+1}-x_k)U_k.
\]
Left, right, and trapezoid sums are finite variants of this operation; their
comparison is rational interval arithmetic, not a limit theorem.

\begin{figure}[!htbp]\centering
\arctanRectangleEnclosureAnimation
\caption{A finite lower/upper rectangle computation.}
\label{fig:arctan-rectangle-enclosure}\end{figure}

\begin{verbatim}
sum := [0,0]
for each cell [x_k,x_(k+1)]:
    sum := sum + width(cell) * certified_range(f, cell, stage)
return sum
\end{verbatim}

\lean{ComputableAnalysis.riemannLeftInterval_add}
\lean{ComputableAnalysis.riemannRightInterval}
\lean{ComputableAnalysis.riemannTrapezoidInterval}
\lean{ComputableAnalysis.riemannTrapezoidInterval_eq_half_add_left_right}
\lean{ComputableAnalysis.riemannTrapezoidInterval_overlap_of_samples}
\leanok

\section{Validity from finite regularity}

A valid integral raw has ordered, nested boxes whose widths eventually fit
every positive rational tolerance. The certificate comes from explicit
regularity data, never from ``all continuous functions are integrable.''

The first reusable theorem accepts a rational Lipschitz bound. On dyadic cells
its width is bounded by \(2L/2^n\):

\begin{verbatim}
rational evaluator + Lipschitz L -> dyadic cell boxes
                                -> width <= 2*L/2^n -> valid raw
\end{verbatim}

\lean{ComputableAnalysis.Integral.LipschitzOnUnit}
\lean{ComputableAnalysis.LipschitzDyadic.compute_nested}
\lean{ComputableAnalysis.LipschitzDyadic.compute_width}
\lean{ComputableAnalysis.LipschitzDyadic.raw_valid}
\lean{ComputableAnalysis.LipschitzDyadic.raw_equiv_of_lipschitz_bounds}
\leanok

For monotone functions, endpoint ranges provide the boxes. A function with
finitely many certified turns is split into monotone pieces. Thus \(|x|\) is
two affine integrals, and \(x|x|/2\) gives its first finite-turn FTC identity.

\lean{ComputableAnalysis.Integral.MonotoneDarbouxSchedule}
\lean{ComputableAnalysis.Integral.monotoneDarbouxScheduleRaw_valid}
\lean{ComputableAnalysis.Integral.PiecewiseMonotoneConstructionFor}
\lean{ComputableAnalysis.Integral.piecewiseMonotoneIntegralFor_equiv_finiteRawSum}
\lean{ComputableAnalysis.Integral.absOnUnit_piecewise_primitiveFTC}
\leanok

\section{The effective FTC}

The central contract is finite telescoping:

\begin{verbatim}
derivative box on a cell -> bound its area -> sum bounds
                        -> telescope primitive increments -> F(b)-F(a)
\end{verbatim}

It is an interface a particular primitive earns. The project does not assert
a universal FTC merely because a function has a classical derivative. Affine
and polynomial witnesses are the regression base; arctangent, exponential,
and the tangent chart are non-polynomial instances.

\lean{ComputableAnalysis.DerivativeBoundFTC}
\lean{ComputableAnalysis.DerivativeBoundFTC.boundedIntegral_equiv_endpointDifference}
\lean{ComputableAnalysis.ExactFunction.affineFTCExact}
\lean{ComputableAnalysis.Integral.squareIntegralEffectiveFTC_integral_equiv_one_third}
\lean{ComputableAnalysis.ArctanEffectiveFTC.arctanEffectiveFTC}
\leanok

Convexity is useful because neighboring secants give finite derivative bounds.
Chapter~\ref{ch:effective-calculus} supplies the reusable certificates.

\section{Circle and sine}

For \(\sin(\pi x)\) on \([0,1/2]\), the ends are rational and every interior
dyadic sample has its own stage-dependent nested-radical box.

\begin{figure}[!htbp]\centering
\intervalSineIntegralStageAnimation
\caption{Dyadic cells and independently refined sine-value boxes.}
\label{fig:interval-sine-integral-stage}\end{figure}

The tangent chart gives a completed Riemann--Stieltjes computation. The
equal-dyadic nested-radical candidate is executable and checked through
stages 0--2; its remaining theorem is the all-depth geometric overlap with
the chart, not a hidden completeness step.

\lean{ComputableAnalysis.SinPiIntegral.dyadicNestedRadicalLeftSum}
\lean{ComputableAnalysis.SinPiIntegral.dyadicNestedRadicalLeftSum_width_le_of_stage}
\lean{ComputableAnalysis.SinPiIntegral.dyadicNestedRadicalLeftSum_zero_overlaps_stieltjes}
\lean{ComputableAnalysis.SinPiIntegral.dyadicNestedRadicalLeftSum_one_overlaps_stieltjes}
\lean{ComputableAnalysis.SinPiIntegral.dyadicNestedRadicalLeftSum_two_overlaps_stieltjes}
\lean{ComputableAnalysis.SinPiIntegral.DyadicTangentWitnessFamily}
\leanok

A trigonometric substitution is one finite parameter partition transported to
its rational-circle image. The Riemann--Stieltjes name emphasizes that image
cells need not be equal.

\section{Finite identities and assembly}

Substitution, integration by parts, and adjacent-interval addition begin as
finite mesh identities. They become integral identities only when both raw
computations have validity and overlap certificates.

\begin{verbatim}
two finite mesh presentations -> reindex / telescope -> overlap edge
\end{verbatim}

\lean{ComputableAnalysis.IntegralIdentities.coordinateTimesArctanForwardTwoStageMonotoneDefiniteIdentity}
\lean{ComputableAnalysis.IntegralIdentities.coordinateTimesArctanForwardTwoStageMonotoneIntegral_equiv_arctanGeom_one}
\lean{ComputableAnalysis.Logarithm.piTriangleLogReciprocalIntegral_equiv_four_arctanGeom_one}
\leanok

\section{Public interface}

To formalize a new integral, provide a rational evaluator; cell ranges or a
modulus; nesting and shrinking; and, when a value matters, an overlap with an
independent computation. For an FTC theorem add a primitive and a finite
derivative-bound telescope. That is the admission rule.
